\begin{scanpage}{0042}{25}
The integrals $v_1^{\pm}$, $v_2^{\pm}$ generally are not expressed in terms of well-known functions. However, by using the method of steepest descents, it is not difficult to obtain their asymptotic expansion when $kr\gg1$. Far from the directions $\phi=\pi\pm\phi_0$ and $\phi=2\alpha-\pi-\phi_0$, the first term of the asymptotic expansion gives us the cylindrical wave diverging from the wedge edge. In the case of wedge excitation by a plane wave (1.23), these cylindrical waves are determined by the equation
\[\left.\begin{aligned}E_z=-H_\phi&=E_{0z}\cdot f^0\cdot\frac{e^{i(kr+\frac\pi4)}}{\sqrt{2\pi kr}},\\E_\phi&=H_z=0,\end{aligned}\right\}\tag{3.31}\]
and with the excitation of the wedge by a plane wave (1.24) they are determined by the equation
\[\left.\begin{aligned}H_z=E_\phi&=H_{0z}\cdot g^0\cdot\frac{e^{i(kr+\frac\pi4)}}{\sqrt{2\pi kr}},\\H_\phi&=E_z=0.\end{aligned}\right\}\tag{3.32}\]
The functions $f^0$ and $g^0$ have the form
\[\left.\begin{aligned}f^0&=\frac{\sin\phi_0}{\cos\phi+\cos\phi_0},\\g^0&=-\frac{\sin\phi}{\cos\phi+\cos\phi_0},\end{aligned}\right\}\tag{3.33}\]
if one face of the wedge $(0<\phi_0<\alpha-\pi)$, and
\[\left.\begin{aligned}f^0&=\frac{\sin\phi_0}{\cos\phi+\cos\phi_0}+\frac{\sin(\alpha-\phi_0)}{\cos(\alpha-\phi)+\cos(\alpha-\phi_0)},\\g^0&=-\frac{\sin\phi}{\cos\phi+\cos\phi_0}-\frac{\sin(\alpha-\phi)}{\cos(\alpha-\phi)+\cos(\alpha-\phi_0)},\end{aligned}\right\}\tag{3.34}\]
if both of its faces $(\alpha-\pi<\phi_0<\pi)$ are illuminated. The index ``0'' for the functions $f^0$ and $g^0$ means that the cylindrical waves (3.31) and (3.32) are radiated by the uniform part of the surface current $(\mathbf j^0)$.
\end{scanpage}
